% Part: normal-modal-logic
% Chapter: filtrations
% Section: preliminaries

\documentclass[../../../include/open-logic-section]{subfiles}

\begin{document}

\olfileid[pt]{nml}{fil}{pre}

\olsection{Preliminares}

As filtrações nos permitem estabelecer a decidibilidade de nossos
sistemas de lógica modal mostrando que eles têm a \emph{propriedade
  dos modelos finitos}, O.s.\@, que toda !!{formula} verdadeira
(falsa) em um modelo também é verdadeira (falsa) em um modelo
\emph{finito}. As filtrações são definidas relativamente a conjuntos
de !!{formula}s fechados sob subfórmulas.

\begin{defn}\ollabel{defn:modallyclosed}
  Um conjunto $\Gamma$ de !!{formula}s é \emph{fechado sob
  subfórmulas} se contém toda subfórmula de !!a{formula} em $\Gamma$.
  Além disso, $\Gamma$ é \emph{modalmente fechado} se é fechado sob
  subfórmulas e, ademais, $!A \in \Gamma$ implica $\Box!A,
  \Diamond!A \in \Gamma$. 
\end{defn}

Por exemplo, dada !!a{formula}~$!A$, o conjunto de todas as suas
sub-!!{formula}s é fechado sob sub-!!{formula}s. Quando definimos uma
filtração de um modelo pelo conjunto de sub-!!{formula}s de~$!A$, ela
tem a propriedade que buscamos: torna $!A$ verdadeira (falsa) se e
somente se o modelo original o faz.

O conjunto de mundos de uma filtração de~$\mModel{M}$ por~$\Gamma$ é
definido como o conjunto de todas as classes de equivalência da
seguinte relação de equivalência.

\begin{defn}
Seja $\mModel{M} =\tuple{W, R, V}$ e suponha que $\Gamma$ seja fechado
sob sub-!!{formula}s. Defina uma relação $\equiv$ em $W$ que valha para
quaisquer dois mundos que tornem verdadeiras as mesmas !!{formula}s de
$\Gamma$, O.s.:
\[
u \equiv v \quad \text{se e somente se } \quad
\forall !A \in \Gamma : \mSat{M}{!A}[u] \Leftrightarrow \mSat{M}{!A}[v].
\]
A classe de equivalência~$[w]_\equiv$ de um mundo~$w$, ou simplesmente
$[w]$, é o conjunto de todos os mundos $\equiv$-equivalentes a~$w$:
\[
[w] = \Setabs{v}{v \equiv w}.
\]
\end{defn}

\begin{prop}
  Dados $\mModel{M}$ e $\Gamma$, a relação $\equiv$ definida acima é
  uma relação de equivalência, O.s.\@, é reflexiva, simétrica e
  transitiva.
\end{prop}

\begin{proof}
  A relação $\equiv$ é reflexiva, pois $w$ torna verdadeiras exatamente
  as mesmas !!{formula}s de~$\Gamma$ que ele próprio. É simétrica,
  pois, se $u$ torna verdadeiras as mesmas !!{formula}s de~$\Gamma$
  que $v$, o mesmo vale para $v$ e~$u$. Também é transitiva, pois, se
  $u$ torna verdadeiras as mesmas !!{formula}s de~$\Gamma$ que~$v$, e
  $v$ as mesmas que $w$, então $u$ torna verdadeiras as mesmas
  !!{formula}s de~$\Gamma$ que~$w$.
\end{proof}

A relação $\equiv$, como qualquer relação de equivalência, divide~$W$
em \emph{partes}, O.s.\@, subconjuntos de~$W$ dois a dois disjuntos
que, em conjunto, cobrem todo~$W$. Todo $w \in W$ é !!a{element} de
uma das partes, a saber, de~$[w]$, pois $w \equiv w$. Portanto, as
partes $[w]$ cobrem todo~$W$. Elas são duas a duas disjuntas, pois, se
$u \in [w]$ e $u \in [v]$, então $u \equiv w$ e $u \equiv v$ e, pela
simetria e transitividade, $w \equiv v$; logo, $[w] = [v]$.

\end{document}
